\chapter{Architecture}\label{ch:architecture}


Due to my limited experience in software development, the project does not focus on performance or software quality.
Therefore the software is written to implement the wanted functionality and should be seen as experimental, since it is neither optimised for performance, nor have proper error handling mechanisms been implemented.

\section{Hardware}

The Project is realized using a RaspberryPI as the central device, connected to which are an external network card, an ESP32 microcontroller, a display and a rotary encoder for user input. 

As the main computer the RaspberryPI 4B with 8GB of RAM was chosen due to its ease of use, linux based operating system and widely available documentation. 

To be able to also test on WPA3 networks, an external network card is needed, due to the internal network card of the RPI not supporting the ciphers used by WPA3.
The Fenvi AX1800 USB WiFi card was chosen for its compatibility with the linux kernel, according to the USB-Wifi github by user morrownr \cite{morrownr_USB_WiFi}.
Another factor for this card was the low pricepoint as well as its smaller formfactor, which lets it fit inside case for the device.

Selecting from the large offer of ESP32 MCUs, again size was taken into account and the ESP32-C6 Mini by Waveshare was chosen for its small size and good feature set.
It connects to the RPI's GPIO pins using simple jumper cables, which supplies the ESP32 with power and establishes a UART connection.

Another device connected to the RPI's GPIO header is the rotary encoder.
On its five pins, it receives 3.3V supply voltage and outputs a gray code, indicating the direction of rotation, on two of its outputs. 
The fith pin is used for the pushbutton inside the rotary encoder.


\begin{figure}[h]
    \centering
    \includegraphics[width=0.9\linewidth]{figures/Abbildungen/Devices.png}
    \label{fig:devices}
    \caption{Hardware Devices}
\end{figure}

\section{Software}

The software of the project is mostly modular and it consists of a central program, which runs the user interface and then executes various shell scripts, with which the functions are implemented.
When executing the setup script, the main program, which is written in C++, gets compiled and the executable gets called by the start.sh script.
The program then displays the user interface, which is a simple interface, printed into the console. 
Each page of the UI is a "Menu" object, each menu has a title, options of the type "MenuOptions", a variable to track the currently selected option, a method to add options to a menu and methods for navigation, input capture and displaying.
MenuOption objects have a name, as well as an action of the std::function type, which can be executed by calling the execute() method.
In the main of the program, the constructors for menues and the addOption() methods are beeing used, to create the different menues and their submenues.
The program starts by calling the navigate() function of the main menu, which then displays all available submenues and menu options.
All submenues work in the same way as the main menu, getting their respective navigate() called upon selection.
During navigation the UI peridiodically prints the available options, using unicode to highlight the currently selected option as well as to give contrast to the title of the menu.

\begin{figure}[h]
    \includegraphics[width=0.3\linewidth]{figures/Abbildungen/hackerypi/menu_main.png}
    \centering
    \caption{Main Menu}
    \label{fig:menues}
\end{figure}

\begin{figure}[h]
    \includegraphics[width=0.25\linewidth]{figures/Abbildungen/hackerypi/menu_wifi.png}
    \includegraphics[width=0.3\linewidth]{figures/Abbildungen/hackerypi/menu_webapp.png}
    \includegraphics[width=0.35\linewidth]{figures/Abbildungen/hackerypi/menu_dev.png}
    \centering
    \caption{Device Menues}
    \label{fig:menues}
\end{figure}
    
For input the device uses a rotary encoder with a pushbutton, turning the encoder moves the selection up or down in the menues, the pushbutton executes the selected option.
The encoder is connected to the RaspberryPIs GPIO interface and the input signals are read and processed by the encoder program.

\lstinputlisting[
    firstline=13,
    lastline=24,
    language=C,
    caption={Rotation handling ISR}
]
    {listings/Raspi-CyberSec-Lab-Project/Program/encoder.h}

The start.sh script starts both programs and pipes the stdout of encoder into the stdin of main.
When an action of the encoder is detected, the according command gets sent to main, the commands beeing the unicodes for the up/down arrows and the enter key on a keyboard.
For example, the User turns the the encoder one click clockwise, therefore main receives one downwards arrow symbol and the selections moves to the next option below.
Because no special commands are utilized for navigation, the program can be easily tested on any computer by navigating with the keyboard.

When an action is executed, main calls the according shell script with the system() function.
The wifiActivate.sh script is used to open wifi hotspots, where the user can choose the type of Wifi standard beeing used. 
It receives an input command and then establishes a WEP, WPA, WPA2 or WPA3 network.

\lstinputlisting[
    language=bash,
    caption={extract of wifiActivate.sh},
    lastline=16,
    breaklines=true,
    basicstyle=\ttfamily\fontsize{8}{9}\selectfont
]
    {listings/Raspi-CyberSec-Lab-Project/Skripte/wifiActivate.sh}
    
With the wifiReset.sh script, the user can deactivate the active network.
Once a network has been cracked, the wifiNewPassword.sh script can be used to give a new random password to the selected network.
The script the pulls a random password out of a password list, which contains popular passwords that have been filtered to work with the specification of each network standard.
To display the current status of the network, meaning its security standard, SSID and password, the wifiStatus.sh script is used.


\begin{figure}[h]
    \includegraphics[width=0.24\linewidth]{figures/Abbildungen/hackerypi/activate_wifi.png}
    \includegraphics[width=0.24\linewidth]{figures/Abbildungen/hackerypi/status_wifi.png}
    \includegraphics[width=0.24\linewidth]{figures/Abbildungen/hackerypi/activate_monitor.png}
    \includegraphics[width=0.24\linewidth]{figures/Abbildungen/hackerypi/menu_password.png}
    \centering
    \caption{WiFi Status and Monitor}
    \label{fig:status_monitor}
\end{figure}

After a WiFi hotspot has been activated, a command with the connection details is sent to the ESP32 over uart.
This functionality is implemented by the uartSendNI.sh script, which gets called from wifiActivate.sh.
Receiving the network information, the ESP32 will then connect to the provided network and act as a communication partner for the RPI.

Selecting the "monitor" option, the user is presented with the options of de-/activating the network monitoring as well as viewing or deleting the monitoring log.
When activating monitoring, wifiMonitor.sh uses the \lstinline[]|iw event -T| command to log devices connecting and disconnecting from the network.

\lstinputlisting[
    language=bash,
    caption={extract of wifiActivate.sh},
    breaklines=true,
    basicstyle=\ttfamily\fontsize{8}{9}\selectfont
]
    {listings/Raspi-CyberSec-Lab-Project/Skripte/uartSendNI.sh}

Inside the webapp menu the JuiceShop and MQTT communication can be found.
When activating the JuiceShop, the juiceShop.sh script will start up the server for the OWASP JuiceShop project and open a hotspot, which can be used to access the simulated webstore.

When a WiFi network has already been activated from the WiFi menu, a simulated MQTT communication can be started from the MQTT menu.
Starting a conversation will again send a command via uart to the ESP32, prompting it to connect to the MQTT broker, running on the RPI.
Following up uart.sh starts the uart program, which starts an infinite MQTT conversation with the ESP32.

The development menu contains options used during development, these can be adjusted on the fly in the development.sh script.

When selecting the shutdown option, all running processes are stopped and the shutdown.sh script calls \lstinline[]|shutdown now| to turn off the device.




